\SourceLeaf{196-erchinger}
\Needspace{8\baselineskip}
\pdfbookmark[1]{Erchinger, seventeen-sided polygon. Review of December 19, 1825}{notice-1825-12-19}
\begin{center}
\rule{.86\linewidth}{.3pt}\par
Göttingische gelehrte Anzeigen.\quad December 19, 1825.\par
\rule{.86\linewidth}{.3pt}
\end{center}
A short paper has been presented to the Royal Society on behalf of Mr. \textsc{Erchinger} of Thuningen in the Kingdom of Würtemberg, on the
\begin{center}\emph{Geometric Construction of the Regular Seventeen-sided Polygon}.\end{center}
As is well known, the general theory of regular polygons has acquired a new form and a wider scope through its close connection with higher arithmetic. The theory of those polygons that can be constructed geometrically is one part of it, though a comparatively small one. Since the time of the Greeks it had been known that the triangle, pentagon, fifteen-sided polygon and all polygons obtained from these by doubling or repeatedly doubling the number of sides have this property, and it was believed, and sometimes expressly asserted, that these were the only ones. Higher arithmetic has shown this to be an error: by laying open the true sources of the entirely general theory, it revealed of itself that, besides the polygons mentioned, there are countless others that can be constructed geometrically, of which the seventeen-sided polygon is the simplest. This shows in the clearest light the superiority of analysis, which encompasses the most general and the particular with equal ease, over geometry, which must always remain with the particular, is held back by ever greater complications as it advances from simpler cases to more complex ones, and would hardly ever have reached, without help from elsewhere, that next case after the known ones. Nevertheless, it remains important, interesting and desirable that purely geometric treatments should also continue to be cultivated, and that
\SourceLeaf{197}
geometry should take possession of at least part of the new fields conquered by analysis. The reviewer knows of no one who has so far publicly treated the construction of the seventeen-sided polygon except Mr. \textsc{Pauker}, in the writings of the Courland Society and in his geometry. Mr. \textsc{Erchinger}'s construction differs from this and is carried out more in the purely geometric spirit. It is as follows. (Anyone can draw the accompanying figure: a straight line on which the points $DBGAIFCE$ lie in that order.) Take any straight line $AB$ and extend it backwards and forwards to $C$ and $D$, so that $AC\times BC=AD\times BD=4AB\times AB$; next determine the points $E$, $G$ on the two sides of the extended line $CA$, so that $AE\times EC=AG\times CG=AB\times AB$, and the point $F$ on the $A$ side of the extended line $BA$, so that $AF\times DF=AB\times AB$; finally divide $AE$ at $I$, so that $AI\times EI=AB\times AF$, where $AI$ is the smaller and $EI$ the larger part of $AE$. Then construct a triangle in which two sides are each $=AB$, and the third $=AI$. If a circle is drawn around this triangle, $AI$ will be the side of the regular seventeen-sided polygon inscribed in that circle.

If one checks the correctness of this construction by comparing it with the theory of the seventeen-sided polygon given as an example in Article 354 of the \emph{Disquisitiones Arithmeticae}, one readily sees that it is simply the geometric translation of the equations to which the application of the general theory leads. Indeed, the distances of the points $C$, $D$, $E$, $F$, $G$, $I$ from $A$ are precisely the quantities denoted there by $(8\pdot1)$, $(8\pdot3)$, $(4\pdot1)$, $(4\pdot3)$, $(4\pdot9)$, $(2\pdot1)$, if positive and negative signs are expressed by position, and the distance of the point $B$ from $A$, taken in the same sense, is set $=-1$. But the real merit of Mr. \textsc{Erchinger}'s paper lies less in setting out the construction itself, since analysis had already marked out the simplest path, than in giving a purely geometric proof of its correctness. This is carried out with such exemplary, painstaking care to avoid everything that is not purely elementary that it does the author credit, and prompts the wish that his indeed uncommon mathematical talent may receive every encouragement.
